\documentclass[addpoints,spanish, 12pt,a4paper]{exam}
%\documentclass[answers, spanish, 12pt,a4paper]{exam}
%\printanswers
\pointpoints{punto}{puntos}
\hpword{Puntos:}
\vpword{Puntos:}
\htword{Total}
\vtword{Total}
\hsword{Resultado:}
\hqword{Ejercicio:}
\vqword{Ejercicio:}

\usepackage[utf8]{inputenc}
\usepackage[spanish]{babel}
\usepackage{eurosym}
%\usepackage[spanish,es-lcroman, es-tabla, es-noshorthands]{babel}


\usepackage[margin=1in]{geometry}
\usepackage{amsmath,amssymb}
\usepackage{multicol}
\usepackage{yhmath}

\usepackage{verbatim}
%\usepackage{pstricks}


\usepackage{graphicx}
\graphicspath{{../img/}} 

\newcommand{\class}{4º Académicas}
\newcommand{\examdate}{\today}
\newcommand{\examnum}{Sistemas de ecuaciones e inecuaciones}
\newcommand{\tipo}{C}


\newcommand{\timelimit}{50 minutos}

\renewcommand{\solutiontitle}{\noindent\textbf{Solución:}\enspace}


\pagestyle{head}
\firstpageheader{\includegraphics[width=0.2\columnwidth]{header_left}}{\textbf{Departamento de Matemáticas\linebreak \class}\linebreak \examnum}{\includegraphics[width=0.1\columnwidth]{header_right}}
\runningheader{\class}{\examnum}{Página \thepage\ of \numpages}
\runningheadrule
\pointsinrightmargin % Para poner las puntuaciones a la derecha. Se puede cambiar. Si se comenta, sale a la izquierda.
\extrawidth{-2.4cm} %Un poquito más de margen por si ponemos textos largos.
\marginpointname{ \emph{\points}}

\begin{document}

\noindent
\begin{tabular*}{\textwidth}{l @{\extracolsep{\fill}} r @{\extracolsep{6pt}} }
\textbf{Nombre:} \makebox[3.5in]{\hrulefill} & \textbf{Fecha:}\makebox[1in]{\hrulefill} \\
 & \\
\textbf{Tiempo: \timelimit} & Tipo: \tipo 
\end{tabular*}
\rule[2ex]{\textwidth}{2pt}
Esta prueba tiene \numquestions\ ejercicios. La puntuación máxima es de \numpoints. 
La nota final de la prueba será la parte proporcional de la puntuación obtenida sobre la puntuación máxima.
\begin{center}


\addpoints
 %\gradetable[h][questions]
	\pointtable[h][questions]
\end{center}

\noindent
\rule[2ex]{\textwidth}{2pt}

\begin{questions}

\begin{comment}
\question[1] 
\begin{solution} \end{solution}

\addpoints
\end{comment}

\question[2] Resuelve el siguiente sistema:
 $$\left\{ \begin{gathered}
  \frac{{x - 4}}{2} + \frac{{x + 2}}{3} \leqslant 2 \hfill \\
  \frac{x}{3} - \frac{x}{2} < 0 \hfill \\ 
\end{gathered}  \right.$$
	\begin{solution} $\left(0, 4\right]$ \end{solution}

\question %La suma de los cuadrados de dos números positivos es 34. Sabiendo que uno es dos unidades mayor que el otro, se pide calcularlos de la siguiente forma:
Pulir el parqué de una habitación ha costado 576\euro. Sabiendo que tiene forma rectangular, que de largo es tres veces su anchura y que el precio del pulido es 12\euro \ el metro cuadrado, calcula las dimensiones de la habitación mediante el siguiente procedimiento:
\begin{parts}
\part[2] Traduce a lenguaje algebraico el enunciado anterior
\begin{solution}  $\left. \begin{gathered}
	  \left(x\cdot y\right)\cdot12=576 \hfill \\
	  x=3y \hfill \\ 
	\end{gathered}  \right\rbrace$ \end{solution}
\part[1] Resuelve la expresión del apartado anterior, indicando cuántas soluciones tiene el problema
\begin{solution}  
%solve([2*x+2*y-34,x**2+y**2-169],[x,y])
$\to 3y\cdot y = \frac{576}{12} \to 3y^2= 48\to\begin{pmatrix}12, & 4\end{pmatrix}$\end{solution}
\end{parts}

\question Resuelve las siguientes inecuaciones:
\begin{parts}

\part[1]$ x^{2} - x - 2 \geq 0  $  
\begin{solution} $ \left(-\infty, -1\right] \cup \left[2, \infty\right)$ \end{solution}

\part[2]$x^{4} - x^{3} + 2 x^{2} - 8 x < 0$  
% solve_univariate_inequality (expand(x*(x-2)*(x-3))<0,x,relational =false )
\begin{solution} $\left(0, 2\right)$ \end{solution}

\part[2]$\dfrac{x^{2} - 1}{x^{3} + 2 x^{2} + x}\geqslant 0$  
%from sympy.solvers.inequalities import reduce_rational_inequalities
%reduce_rational_inequalities([[(2*x-2)/(1-3*x) < -2/3]], x,relational=0)
%reduce_rational_inequalities([[expand(x*(x-1))/expand((x+1)**1) >= 0]], x,relational=0)
%reduce_rational_inequalities([[expand((x-1)*(x+1))/expand(x*(x+1)**2) >= 0]], x,relational=0)
\begin{solution} $\dfrac{\left(x-1\right)\left(x+1\right)}{x\left(x+1\right)^2}\to\left(-1, 0\right) \cup \left[1, \infty\right)$ \end{solution}

\part[2]  $\left| {2x - 12} \right| > 2$ 
%reduce_abs_inequality(Abs(2*x - 12) - 2, '>', x)
\begin{solution} $\left(-\infty, 5\right) \cup \left(7, \infty\right) $ \end{solution}

\end{parts}

\addpoints
\question La tarifa de  telefonía de la empresa A es 20 \euro\ fijos mensuales más 7 céntimos de euro por minuto de conversación, la de la empresa B es 11 \euro\ fijos más 12 céntimos por minuto de conversación. ¿A partir de cuantos minutos empieza a ser más rentable la tarifa de la empresa A? Resuelve el problema de la siguiente forma
\begin{parts} 
\part[2] Traduce a lenguaje algebraico el enunciado anterior
\begin{solution} $\left. \begin{gathered}
	  20+0,07x<11+0,12x \hfill \\ 
	 \end{gathered}  \right\rbrace $ \end{solution}
\part[1] Resuelve la expresión del apartado anterior e indica cuáles son las soluciones
\begin{solution} $\left(180.0, \infty\right) \to $ A partir de 180 m o 3 horas  \end{solution}

\end{parts}
\addpoints

\end{questions}

\end{document}
\grid
